% Part:sets-functions-relations
% Chapter: size-of-sets
% Section: non-enumerability

\documentclass[../../../include/open-logic-section]{subfiles}

\begin{document}

\olfileid{sfr}{siz}{nen}

\olsection{不可数集}

\begin{editorial}
  本节使用 \olref[enm]{sec} 中的定义来证明 $\Bin^\omega$ 和
  $\Pow{\PosInt}$ 的不可数性。相比 \olref[enm-alt]{sec} 中的版本，
  此版本设计得更为基础，也更为详细。
\end{editorial}

有些集合，比如正整数集合 $\PosInt$，是无限的。
到目前为止，我们见过的无限集合都是可数的。
然而，也存在不具备这种性质的无限集合。
这样的集合称为\emph{不可数的}。

首先，不可数集合的存在或许已经令人惊讶。
对于任何可数集合~$A$，都存在一个满射函数 $f \colon \PosInt \to A$。
如果一个集合是不可数的，则不存在这样的函数。
也就是说，任何将~$\PosInt$ 中无限多个元素映射到~$A$ 的函数都无法穷尽~$A$ 的所有元素。
因此，~$A$ 的元素“多于”无限多的正整数。

那么，如何证明一个集合是不可数的呢？
你必须证明这样的满射函数不可能存在。
等价地，你必须证明~$A$ 的元素不能被枚举为一个单向的无限列表。
最好的办法是证明，~$A$ 的任何元素列表都必然至少漏掉一个元素；
或者说，任何函数 $f\colon \PosInt \to A$ 都不可能是满射。
我们可以使用 Cantor 的\emph{对角方法}来做到这一点。
给定~$A$ 的一个元素列表，比如 $x_1$, $x_2$, \dots，
我们构造出~$A$ 的另一个元素，根据其构造方式，它不可能出现在该列表中。

我们的第一个例子是所有由 $0$ 和 $1$ 组成的无限且无间隔的序列的集合~$\Bin^\omega$。

\begin{thm}
\ollabel{thm:nonenum-bin-omega}
$\Bin^\omega$~是不可数的。
\end{thm}

\begin{proof}
为导出矛盾，假设 $\Bin^\omega$ 是可数的，即假设存在一个列表 $s_{1}$, $s_{2}$, $s_{3}$, $s_{4}$, \dots{}，包含了~$\Bin^\omega$ 的所有元素。
每个 $s_i$ 本身都是一个由 $0$ 和~$1$ 组成的无限序列。
让我们称此列表中第 $i$ 个序列的第 $j$ 个元素为 $s_i(j)$。
那么第 $i$ 个序列~$s_i$ 是
\[
s_i(1), s_i(2), s_i(3), \dots
\]

我们可以将这个列表，以及其中每个序列 $s_i$ 的元素，排列成一个阵列：
\[
\begin{array}{c|c|c|c|c|c}
& 1 & 2 & 3 & 4 & \dots \\\hline
1 & \mathbf{s_{1}(1)} & s_{1}(2) & s_{1}(3) & s_1(4) & \dots \\\hline
2 & s_{2}(1)& \mathbf{s_{2}(2)} & s_2(3) & s_2(4) & \dots \\\hline
3 & s_{3}(1)& s_{3}(2) & \mathbf{s_3(3)} & s_3(4) & \dots \\\hline
4 & s_{4}(1)& s_{4}(2) & s_4(3) & \mathbf{s_4(4)} & \dots \\\hline
\vdots & \vdots & \vdots & \vdots & \vdots & \mathbf{\ddots}
\end{array}
\]
侧边的标签给出了列表中序列的编号 $s_1$, $s_2$, \dots；顶部的数字则标明了各个序列的元素序号。
例如，$s_{1}(1)$ 表示序列 $s_{1}$ 的第一个元素是哪个数字（$0$ 或~$1$），依此类推。

现在我们来构造一个由 $0$ 和 $1$ 组成的无限序列 $\overline{s}$，它不可能出现在这个列表中。
$\overline{s}$ 的定义将依赖于列表 $s_1$, $s_2$, \dots。
任意一个由 0 和 1 的无限序列组成的无限列表，都会产生一个保证不在该列表中的无限序列~$\overline{s}$。

为了定义 $\overline{s}$，我们指定其所有元素是什么，即对所有的 $n \in \PosInt$ 指定 $\overline{s}(n)$。
我们通过沿上面阵列的对角线读取（因此得名“对角线方法”），并将每一个 $1$ 改为 $0$、每一个 $0$ 改为~$1$ 来完成这一点。
更抽象地说，我们根据对角线上的第 $n$ 个元素 $s_n(n)$ 是 $1$ 还是 $0$，将 $\overline{s}(n)$ 定义为 $0$ 或 $1$。
\[
\overline{s}(n) =
\begin{cases}
1 & \text{若 $s_{n}(n) = 0$}\\
0 & \text{若 $s_{n}(n) = 1$}.
\end{cases}
\]
如果你更喜欢公式而不是分情况定义，也可以定义 $\overline{s}(n) = 1 - s_n(n)$。

显然 $\overline{s}$ 是一个由 $0$ 和 $1$ 组成的无限序列，因为它不过是阵列对角线上 $0$ 和 $1$ 序列的“镜像”序列。
所以 $\overline{s}$ 是~$\Bin^\omega$ 的一个元素。
但它不可能在列表 $s_1$, $s_2$, \dots{} 中。为什么呢？

它不可能是列表中的第一个序列 $s_1$，因为它的第一个元素与 $s_1$ 不同。
无论 $s_1(1)$ 是什么，我们都将 $\overline{s}(1)$ 定义为取反后的值。
它也不可能是列表中的第二个序列，因为 $\overline{s}$ 在第二个元素上与 $s_2$ 不同：如果 $s_2(2)$ 是 $0$，则 $\overline{s}(2)$ 是 $1$，反之亦然。
依此类推。

更精确地说：如果 $\overline{s}$ 在列表中，就会存在某个 $k$ 使得 $\overline{s} = s_{k}$。
两个序列相同当且仅当它们在每个位置上都一致，即对任意~$n$，$\overline{s}(n) =
s_{k}(n)$。
因此特别地，取 $n = k$ 这一特殊情况，$\overline{s}(k) = s_{k}(k)$ 必须成立。
$s_k(k)$ 要么是 $0$，要么是~$1$。
如果它是 $0$，那么 $\overline{s}(k)$ 必须是~$1$——我们正是这样定义 $\overline{s}$ 的。
但如果 $s_k(k) = 1$，那么同样根据 $\overline{s}$ 的定义，$\overline{s}(k) = 0$。
无论哪种情况，$\overline{s}(k) \neq s_{k}(k)$。

我们一开始假设存在一个 $\Bin^\omega$ 的元素列表 $s_1$, $s_2$, \dots{}
从这个列表我们构造了一个序列~$\overline{s}$，并证明了它不可能在该列表中。
然而，如果所有的 $s_i$ 都是 $0$ 和 $1$ 的序列，那么 $\overline{s}$ 肯定也是一个 $0$ 和 $1$ 的序列，即 $\overline{s} \in \Bin^\omega$。
这尤其表明，不可能存在一个包含~$\Bin^\omega$ \emph{所有}元素的列表，因为对于任何这样的列表，我们都能构造出一个保证不在其上的序列~$\overline{s}$，因此，假设存在一个包含~$\Bin^\omega$ 中所有序列的列表会导致矛盾。
\end{proof}

\begin{explain}
这种证明方法之所以称为“对角线法”，是因为它利用阵列的对角线来定义~$\overline{s}$。
对角线法并不一定需要阵列的存在：即使没有实际的阵列和对角线，我们也可以使用类似的思想来证明集合是不可数的。
\end{explain}

\begin{thm}
\ollabel{thm:nonenum-pownat}
$\Pow{\PosInt}$ 是不可数的。
\end{thm}

\begin{proof}
我们采用同样的方法，证明对于~$\PosInt$ 的任意子集列表，总存在一个不在该列表中的 $\PosInt$ 的子集。
假设给定了~$\PosInt$ 的如下子集列表：
\[
Z_{1}, Z_{2}, Z_{3}, \dots
\]
现在定义一个集合 $\overline{Z}$，使得对任意 $n \in \PosInt$，$n \in \overline{Z}$ 当且仅当 $n \notin Z_{n}$：
\[
\overline{Z} = \Setabs{n \in \PosInt}{n \notin Z_n}
\]
显然 $\overline{Z}$ 是一个正整数集合，因为根据假设每个~$Z_n$ 都是正整数集合，因此 $\overline{Z} \in \Pow{\PosInt}$。但 $\overline{Z}$ 不可能在该列表中。为此，我们将证明对于每个 $k \in \PosInt$，都有 $\overline{Z} \neq Z_k$。

任取 $k \in \PosInt$。我们定义 $\overline{Z}$ 使得对任意 $n \in \PosInt$，$n \in \overline{Z}$ 当且仅当 $n \notin Z_n$。
特别地，取 $n=k$，则有 $k \in \overline{Z}$ 当且仅当 $k \notin Z_k$。
但这表明 $\overline{Z} \neq Z_k$，因为 $k$ 属于其中一个集合而不属于另一个，所以 $\overline{Z}$ 和 $Z_k$ 具有不同的元素。由于 $k$ 是任意的，故 $\overline{Z}$ 不在列表 $Z_1$, $Z_2$, \dots 中。
\end{proof}

\begin{explain}
前面的证明没有提到对角线，但你可以这样来理解它涉及一条对角线：想象将集合 $Z_1$, $Z_2$, \dots 写成一个数组，其中每个元素~$j \in Z_i$ 被列在第~$j$ 列。假设该列表的前四个集合是 $\{1,2,3,\dots\}$, $\{2, 4, 6, \dots\}$, $\{1,2,5\}$ 和 $\{3,4,5,\dots\}$。那么这个数组的开头将是
\[
\begin{array}{r@{}rrrrrrr}
  Z_1 = \{ & \mathbf{1}, & 2, & 3, & 4, & 5, & 6, & \dots\}\\
  Z_2 = \{ &  & \mathbf{2}, &  & 4, &  & 6, & \dots\}\\
  Z_3 = \{ & 1, & 2, &  &  & 5\phantom{,} &  & \}\\
  Z_4 = \{ &  &  & 3, & \mathbf{4}, & 5, & 6, & \dots\}\\
  \vdots & & & & & \ddots
\end{array}
\]
那么 $\overline{Z}$ 就是通过沿对角线向下走所得到的集合，排除对角线上出现的任何数字，并包含那些在数组的第 $j$ 行/列有空缺的 $j$。在上面的例子中，我们会排除 $1$ 和 $2$，包含~$3$，排除~$4$，等等。
\end{explain}

\begin{prob}
用对角线论证证明 $\Pow{\Nat}$ 是不可数的。
\end{prob}

\begin{prob}\label{sfr:siz:nen:prob:f-posint}
通过显式的对角线论证，证明由函数 $f \colon \PosInt \to \PosInt$ 构成的集合是不可数的。也就是说，证明如果 $f_1$, $f_2$, \dots 是一个函数列表且每个 $f_i\colon \PosInt \to \PosInt$，那么存在某个 $\overline{f}\colon \PosInt \to \PosInt$ 不在此列表中。
\end{prob}

\end{document}